\section{Tag 6 — Polynome über GF($2^3$) und die Reed-Solomon-Idee}

\subsection*{Bleistiftübung 1 — Polynome rechnen}

(a) $f(x) + g(x)$: koeffizientenweise XOR.
\begin{align*}
  &\text{Konstante: } 4 \oplus 5 = 1 \\
  &\text{Koeff. von } x: 3 \oplus 6 = 5 \\
  &\text{Koeff. von } x^2: 1 \oplus 0 = 1 \\
  \Rightarrow\; & f(x) + g(x) = 1 + 5x + x^2.
\end{align*}

(b) $(1+x)^2$ über GF($2^3$):
\begin{align*}
  (1 + x)(1 + x)
  &= 1 + x + x + x^2 \\
  &= 1 + (1 \oplus 1) x + x^2 \\
  &= 1 + 0 \cdot x + x^2 = 1 + x^2.
\end{align*}

\textbf{Vergleich:} in der gewohnten Algebra wäre $(1+x)^2 = 1 + 2x +
x^2$. In GF($2^3$) ist aber $1 + 1 = 0$, der mittlere Term
verschwindet also. Allgemein gilt in jedem Körper mit
\emph{Charakteristik 2} (wie GF($2^n$)) der „Freshman's Dream":
$(a + b)^2 = a^2 + b^2$. Einer der Gründe, warum GF($2^n$) für
Computer-Implementierungen so charmant ist.

(c) $(2 + x) \cdot (3 + x^2)$:
\begin{align*}
  &= 2 \cdot 3 + 2 \cdot x^2 + x \cdot 3 + x \cdot x^2 \\
  &= 6 + 3x + 2x^2 + x^3,
\end{align*}
denn $2 \cdot 3 = 6$ (aus der Multiplikationstafel von Tag 5:
$x \cdot (x+1) = x^2 + x = 6$).

In Standardform: $\mathbf{6 + 3x + 2x^2 + x^3}$.

\subsection*{Bleistiftübung 2 — Gerade durch zwei Punkte}

(a) Aus $f(x) = a_0 + a_1 x$ und den beiden Auswertungen:
\begin{align*}
  a_0 + a_1 \cdot 1 &= 4 \\
  a_0 + a_1 \cdot 2 &= 7
\end{align*}

(b) Subtrahieren wir die erste Gleichung von der zweiten (in GF: das
ist Addition):
\[
  (a_0 + 2 a_1) + (a_0 + a_1) = 7 + 4
  \;\Rightarrow\; (1 \oplus 2) a_1 = 7 \oplus 4
  \;\Rightarrow\; 3 a_1 = 3.
\]
Also $a_1 = 3 \cdot 3^{-1}$. Aus Tag 5 wissen wir: das Inverse von
$3$ in GF($2^3$) ist $6$ (denn $3 \cdot 6 = 1$). Also
$a_1 = 3 \cdot 6 = 1$.

Eingesetzt in die erste Gleichung: $a_0 = 4 + a_1 = 4 + 1 = 5$.

(c) Das Polynom ist $\mathbf{f(x) = 5 + x}$. Verifikation:
\begin{itemize}
\item $f(1) = 5 + 1 = 5 \oplus 1 = 4$ \checkmark
\item $f(2) = 5 + 1 \cdot 2 = 5 \oplus 2 = 7$ \checkmark
\end{itemize}

\subsection*{Bleistiftübung 3 — RS$(5, 3)$ von Hand}

(a) $f(x) = 3 + 5x + 2x^2$, ausgewertet an den fünf Stützstellen:

\begin{center}
\begin{tabular}{cl}
  \toprule
  $x$ & $f(x)$ \\
  \midrule
  $1$ & $3 + 5 + 2 = 4$ \\
  $2$ & $3 + (5 \cdot 2) + (2 \cdot 4) = 3 + 1 + 3 = 1$ \\
  $3$ & $3 + (5 \cdot 3) + (2 \cdot 3^2) = 3 + 4 + (2 \cdot 5) = 3 + 4 + 1 = 6$ \\
  $4$ & $3 + (5 \cdot 4) + (2 \cdot 4^2) = 3 + 2 + (2 \cdot 6) = 3 + 2 + 7 = 6$ \\
  $5$ & $3 + (5 \cdot 5) + (2 \cdot 5^2) = 3 + 7 + (2 \cdot 7) = 3 + 7 + 5 = 1$ \\
  \bottomrule
\end{tabular}
\end{center}

Hilfsrechnungen aus der GF($2^3$)-Multiplikationstafel:
$5 \cdot 2 = 1$, $2 \cdot 4 = 3$, $5 \cdot 3 = 4$, $3 \cdot 3 = 5$,
$2 \cdot 5 = 1$, $5 \cdot 4 = 2$, $4 \cdot 4 = 6$, $2 \cdot 6 = 7$,
$5 \cdot 5 = 7$, $2 \cdot 7 = 5$.

Codewort: $\mathbf{[4, 1, 6, 6, 1]}$.

(b) Wenn $f(2)$ und $f(4)$ als verloren markiert sind und der
Empfänger $f(1) = 4$, $f(3) = 6$, $f(5) = 1$ kennt:

\begin{itemize}
\item Er weiß: das gesuchte Polynom hat Grad $< 3$, also Form
  $a_0 + a_1 x + a_2 x^2$ mit drei unbekannten Koeffizienten.
\item Drei Auswertungen sind drei Gleichungen in den drei
  Unbekannten — ein lineares Gleichungssystem über GF($2^3$).
\item Das System hat genau eine Lösung (aus dem Theorem über die
  Stützstellen-Eindeutigkeit), die er ausrechnen kann. Die
  Datensymbole sind dann genau die Koeffizienten.
\end{itemize}

(c) Bei RS$(5, 3)$ ist $n - k = 2$. Korrigierbar sind also
$\lfloor 2/2 \rfloor = \mathbf{1}$ nicht-lokalisierter Fehler. (Zwei
Auslöschungen sind möglich, aber nur ein „echter" Fehler.)

\subsection*{Aufgaben 1.1, 1.2, 2.1, 2.2 (Python)}

\textbf{1.1 — Horner-Schema:}
\[
  a_0 + a_1 x + a_2 x^2 = a_0 + x \cdot (a_1 + x \cdot a_2).
\]
Statt drei Multiplikationen ($a_1 \cdot x$, $a_2 \cdot x^2 = a_2 \cdot
x \cdot x$) und einer Addition pro Term reicht eine Multiplikation
und eine Addition pro Koeffizient — von außen nach innen geklammert.
Horner ist die effizienteste Polynom-Auswertung.

\textbf{1.2:} bei richtiger Implementierung soll
\texttt{f.auswerten(GF8(2))} genau \texttt{GF8(1)} ergeben (siehe
Bleistiftübung 3). Falls nicht: in der Multiplikationstafel
nachschlagen, wo der Vorzeichenfehler steckt.

\textbf{2.1:} mit derselben Stützstelle zweimal verschwindet die
Eindeutigkeit — zwei gleiche Auswertungen liefern keine zwei
Gleichungen, sondern nur eine. Der Empfänger hätte effektiv ein
Codewort der Länge $n-1$, würde aber Längen-$n$ erwarten. Schlecht.

\textbf{2.2:} Wenn der Empfänger gar nicht weiß, dass
\texttt{codewort[2]} verfälscht wurde, ist es ein „echter" Fehler.
Bei RS$(5, 3)$ kann genau ein solcher korrigiert werden — Tag 8 zeigt
das Verfahren. Wüsste der Empfänger, dass diese Position kaputt ist
(Position-Wissen aus dem Übertragungsmedium), wäre es nur eine
Auslöschung — und davon kann er bis zu zwei korrigieren.
